% !TEX root = ../main_zh.tex
\chapter{导论：宏观经济学的思维方式}
\label{ch:intro}

学完中级微观经济学，一个学生已经能颇为得心应手地分析单个家庭、单个企业或单个市场。消费者把价格和收入当作给定，去挑选自己最偏好的消费组合；企业把工资率和资本租金率当作给定，去决定生产多少。全部的分析功夫都花在\emph{选择}本身上，而进入选择问题的那些数字——这个价格、那笔收入——都来自外部，已经被人测量好了。宏观经济学恰恰从这种便利结束的地方开始。它研究的对象不是某个家庭或某家企业，而是\emph{作为整体}的经济：一个国家的总产出、所有价格的平均水平、愿意工作却找不到工作的人所占的比例、一代人的生活水准提高的速度。没有人会把这些数字递到我们手上。在建立任何一个总量经济模型之前，我们必须先学会如何\emph{度量}它。

这是宏观经济学的第一个特殊之处，也是本书在任何模型出场之前先用三章讨论度量的原因。但还有第二个、更深层的差别，它塑造了整个学科。整体的行为并不是各个部分行为的简单加总。一个家庭决定多储蓄，会提高这个家庭的财富；可如果所有家庭同时都想多储蓄，总支出就会下降，企业卖得更少，收入随之减少，到头来整个经济储蓄的可能并不比从前多。\emph{加总不等于相加}。要对整体说出任何靠得住的话，我们就需要理论——而正如我们将看到的，宏观经济学家用两种截然不同的风格构建了这套理论，一种看重逻辑的纯粹，另一种看重实用。本章先把全书要回答的问题、要用的方法以及它所依托的两大思想传统摆出来。它是后续一切内容的地图。

\section{宏观经济学研究什么}

宏观经济学研究的是作为整体的经济，而不是它的各个组成部分。研究对象一变，我们首先要解决的问题也跟着变了。在微观经济学里，度量相对容易：消费者把收入和价格当作外生给定，企业要核算自己的成本，但牵涉到的那些数量是具体的、有边界的。一旦我们上升到整个国家的层面，连最基本的数量也得被\emph{构造}出来。我们并不像一个购物者那样直接观察到``价格''；我们观察到的是数以百万计的单个价格，必须设法把它们合成一个单一的\emph{价格水平}（第~\ref{ch:prices}~章）。我们也观察不到``产出''；我们观察到的是无数企业以彼此不可通约的单位进行的生产，必须把它们加总成像国内生产总值这样的一个度量（第~\ref{ch:accounts}~章）。

\state{宏观经济学：把经济作为整体来研究}{
微观经济学分析单个的经济主体和市场，把价格与收入当作\emph{给定}。宏观经济学分析的是总量——总产出、价格水平、失业率、长期增长率——这些量并非给定，而是必须先被\emph{度量并加总}出来。这正是本书先讲国民收入核算（第~\ref{ch:accounts}~章）、价格指数（第~\ref{ch:prices}~章）和劳动力市场（第~\ref{ch:labor}~章），然后才转向任何理论的原因。
}

正因为这些总量是构造出来的、而不是直接读取来的，它们在某种程度上也是\emph{被构造的数字}——而被构造的数字会受到构造时所做选择的影响。这与凭空捏造不是一回事。像 GDP 这样的数据，是依照一套确定的统计程序、从真实的底层数据组装出来的，但这套程序包含了主观判断，而这些判断是可以被拿捏的。因此，认真的宏观经济数据读者会拿头条总量去对照那些更难做手脚的物理替代指标，也去对照其它经济体的同类序列。

\rmkb[与官方数字相互印证]{由于头条总量建立在构造选择之上，分析者常常用难以粉饰的物理指标来做交叉验证。一个有名的例子是所谓的克强指数，它不单看上报的产出数字，而是用电力消费量、铁路货运量和银行贷款来衡量中国的经济活动；其前提是，一家工厂烧掉的电量，比它声称创造的产值更难造假。把同一序列拿到不同经济体之间去比较——比如中国对美国——是又一道理智的核验。这一切并不意味着官方数字是错的；它意味着一个严肃的宏观经济学家会把每一个总量都当作一项背后有方法的度量，并追问那套方法可能漏掉了什么。}

\section{宏观经济学家怎样工作}

面对总量经济，人究竟是怎样做宏观经济学的？这门学科有它特有的推进方式。它先是收集\emph{程式化事实}（stylized facts）——数据中那些宽泛而稳健的规律性，比如观察到消费与收入紧密同向变动，或者产出与就业在经济周期中同涨同落——然后试图解释它们。解释可以通过几种不同的方法来进行，值得在一开始就把它们区分清楚，因为本书后面会把它们全用上。

\begin{itemize}
    \item \textbf{简约式分析（reduced-form analysis）。}人们跑回归，把各总量变量联系起来，报告数据中发现的相关关系。这是证据，而且是有用的证据，但它有两个众所周知的局限。一个回归系数未必揭示背后的\emph{机制}；它顶多记录下一种关系。而且宏观经济是一个\emph{动态、环环相扣的系统}，变量之间彼此反馈，所以单个回归看似讲出的那种干净的单向因果故事，往往是站不住脚的。
    \item \textbf{建模。}人们写下一个明确的模型——关于经济主体、技术和市场的若干假设——并推导出它的含义。回报是得到一个由既定前提合乎逻辑地推出的\emph{预测结果}，因此任何预测的来源都可以被追溯，模型也可以被追问。
    \item \textbf{模拟。}人们把历史数据喂进一个模型，看它的行为与实际发生的情况吻合得有多好，借此检验理论与历史记录的拟合程度。
    \item \textbf{校准（calibration）。}人们不是把每个参数都从头重新估计，而是依靠积累起来的经验估计——成堆的既有研究和数据——来确定一个模型的参数，然后再问这样一个被既有经验证据约束住的模型意味着什么。
\end{itemize}

这些方法是互补的，而不是相互竞争的。一个程式化事实激发出一个模型；模型用既有的估计来校准；校准过的模型拿去与历史做模拟；简约式回归既提供了最初的那个事实，又对模型的预测做了核验。读者几乎会在每一章里都看到这个循环以缩微的形式重演。

\rmkb[关于资料与教材]{本课程以新古典传统讲授，英文方面的天然搭档是 Williamson 的《宏观经济学》（\emph{Macroeconomics}），它从明确的微观基础出发展开整个学科。在制度层面——也就是宏观经济政策在中国究竟是怎样运作的——讲授参考了把宏观经济理论与中国政策相搭配的材料。在数据方面，标准的源头是中国国家统计局、CEIC 和 Wind 这类商业宏观数据库，以及在国际和美国序列上常用的美联储经济数据（FRED）服务。知道数字从哪里来，是知道它们意味着什么的一部分。}

\section{两大传统：新古典与凯恩斯}

宏观经济学里有一个根本性的方法论岔路口，理解了它，就能在很大程度上解释这门学科为什么是今天这个样子。两个分支在一个问题上分道扬镳：人应当从个体出发往上构建，还是直接从总量入手？

\subsection{新古典传统：微观基础}

新古典的观点在微观经济学和宏观经济学之间不划出截然的界线。在这幅图景里，原则上根本就没有什么独有的``宏观''现象在发生：经济就是它那些做最优化的消费者和企业的加总，因此宏观经济结果应当可以从个体的微观经济行为中推导出来，并与之相一致。现代新古典宏观经济学于是坚持要有\emph{微观基础}（microfoundations）——每一个总量关系都必须扎根于在约束下最大化效用或利润的经济主体的选择之中。

这种纪律的吸引力是实实在在的。一个有微观基础的模型在\emph{逻辑上是严丝合缝的}：它的结论由关于偏好、技术和约束的明确假设严格推出，因此人们总能确切地说出某个结果为什么成立，并把它一直追溯到一个可供争辩的前提。代价同样是实实在在的。人类社会包含着太多互相牵动的部分，而一个干净、内部自洽的逻辑体系往往很难把它们全部容纳进来。坚持要把每一个宏观经济现象都从第一性原理推出，会让最终得到的模型对于我们实际观察到的那个杂乱经济\emph{解释力}弱得出人意料——任何曾经试图把一个漂亮模型套到顽固数据上的人，都熟悉这种沮丧。

\subsection{凯恩斯传统：让数据说话}

凯恩斯式的回应是改变出发点。与其从个体最优化往上构建总量，不如\emph{直接}研究总量，从总量数据中读出整体的行为，让数据说话——哪怕代价是把微观基础悬置不谈。

最典型的例子是总量消费函数。人们不去从每个家庭的效用最大化推导一个社会消费多少，而是看总量消费数据，提出一个简单的关系，
\[
C = C_0 + \mathrm{MPC}\cdot(Y - T),
\]
其中 $Y - T$ 是总量可支配收入（产出扣去税收），$C_0$ 是自发消费，$\mathrm{MPC}$ 是\emph{边际消费倾向}（marginal propensity to consume）——可支配收入每增加一元中被花掉的那一部分。值得注意的一点是，$\mathrm{MPC}$ 没有任何微观经济学上的来历。它不是从任何偏好结构推导出来的；它只是一个拟合数据的参数，因为数据大致就是那样表现的。同样的精神后面会出现在索罗模型（第~\ref{ch:solow}~章）里，那里直接\emph{假定}一个不变的储蓄率，而不是从一个跨期最优化问题推导储蓄。

\state{凯恩斯式的取舍}{
凯恩斯式的进路用\emph{科学的纯粹性换取实用性}。通过从总量出发、把 $C = C_0 + \mathrm{MPC}\cdot(Y-T)$ 这样的关系直接拟合到数据上，它放弃了对这些关系的微观经济推导。它换来的，是一套可操作、有经验依据的论述，目标直指\emph{理解数据本身}——这样的模型也许能很好地解释世界，尽管按其构造方式，它讲不出每个经济主体为什么会做出总量所暗示的那种行为。
}

两个传统都不是简单地对或错。新古典模型更严谨，将占据我们增长理论的大部分篇幅；凯恩斯模型对短期政策更立竿见影地有用，将组织起本书的后半部分。对宏观经济学成熟的解读，是把两者同时放在心里。

\rmkb[远在凯恩斯之前，人们就在研究总量]{若以为宏观经济问题始于凯恩斯，那就错了。古典经济学家早在他之前就把国民财富作为整体来研究——亚当·斯密的《国富论》在真切的意义上就是一部关于总量经济的著作——他们用的是那个时代的古典方法。凯恩斯改变的不是研究的主题，而是研究的方法：他让人们可以正大光明地按总量自身的逻辑去给它建模，而不必先把它还原为个体。}

\section{两个大问题}

尽管细节散落各处，宏观经济学其实围绕两个问题来组织，二者以时间跨度相区分。

\begin{itemize}
    \item \textbf{长期经济增长。}为什么有些国家富、有些国家穷，又是什么决定了一个经济的产出在数十年里的长期\emph{趋势}？这是增长理论的领地——索罗模型及其有微观基础的后继者（第~\ref{ch:solow}--\ref{ch:malthus}~章）。
    \item \textbf{短期经济周期。}产出为什么不是沿着趋势平滑增长，而是绕着趋势\emph{波动}，有繁荣也有衰退，政策又能不能把这些波动熨平？这是短期宏观经济学的领地——凯恩斯交叉、IS--LM 以及总需求与总供给（第~\ref{ch:keynes}--\ref{ch:adas}~章）。
\end{itemize}

把这两个问题放在一起把握，一个简洁的办法是把产出的路径看成趋势加上偏离：
\[
\underbrace{\text{观察到的产出}}_{\text{我们看到的}}
\;=\;
\underbrace{\text{长期增长趋势}}_{\text{第 \ref{ch:solow}--\ref{ch:malthus} 章}}
\;+\;
\underbrace{\text{绕趋势的周期性冲击}}_{\text{第 \ref{ch:keynes}--\ref{ch:adas} 章}}.
\]
增长理论研究趋势；经济周期理论研究对趋势的偏离。本书正是按这个划分推进的：先度量各总量，再解释趋势，最后解释波动。

\section{宏观经济政策的目标}

这一切对政策又为什么重要？宏观经济政策按惯例瞄准四个目标，而本书余下的部分都可以读作一项研究：研究是什么挡在了实现这些目标的路上。

\begin{enumerate}
    \item \textbf{经济增长}——在长期内提高产出和生活水准。
    \item \textbf{充分就业}——确保每一个主动想工作的人都能找到工作。
    \item \textbf{控制通胀}——把价格水平稳定到足以让货币保持为一把可靠的标尺。
    \item \textbf{国际收支平衡}——保持国际收支的平衡，这是经济一旦对贸易和资本流动开放就会浮现的关切。
\end{enumerate}

有三点澄清值得现在就标出来，每一点后面都有专门的章节展开。

第一，``充分就业''并不意味着失业率为零。它意味着每一个主动找工作的人都能找到工作，从而失业的\emph{周期性}成分为零，测得的失业率只反映\emph{自然失业率}（natural rate）——也就是那种即便在景气时期也会持续存在的摩擦性和结构性失业，源于人们在不同工作和行业之间流动。我们将在劳动力市场一章（第~\ref{ch:labor}~章）把这一点讲精确。

第二，这些目标无法同时达成。当波动由需求驱动时，就业和通胀往往同向变动，于是朝着更充分就业去推，往往会抬高通胀，反之亦然：两者不能同时被最优化。这种张力是菲利普斯曲线（Phillips curve，第~\ref{ch:phillips}~章）的主题，也是稳定政策的核心两难。

\rmkb[中国式的国际收支平衡]{国际收支平衡这个目标，会因一国在世界经济中所处的位置不同而呈现出不同的样貌。中国凭借强劲的出口长期保持着巨额的\emph{贸易顺差}，央行对汇率握有近乎绝对的掌控权。讲授中反复提出的一个问题是：由此积累起来的对外债权，是否真的像它表面看上去那么值钱？堆积对外顺差，在某种意义上就是把你自己的产出借给世界其余部分，而这笔``储蓄''的价值，取决于它最终被换成了什么。本书只在边角处才回到开放经济的问题，但值得早早指出：``国际收支平衡''是一个鲜活而有争议的政策目标，而不是一条已成定论的会计恒等式。}

\section{关于本学科性情的一句话}

一开始就坦诚地说清这是一门什么样的学科，是值得的。微观经济学是齐整的：一个提法得当的问题，核心通常是一个干净的最优化，并有一个确定的答案。宏观经济学则刻意地不那么齐整。同一个事实往往可以透过不止一个模型来解读；数据很少能干净利落地了结一场争论；而一堂课上最富启发的话，有时恰恰是一句一行字的制度性旁白——一家央行实际上是怎样进行某项操作的，某项政策为什么偏偏在某一年失败了，某个国家是怎样度量某个特定序列的——而这样的话哪一章都安放得不算妥帖。我们不会假装并非如此，也不会硬把每一条这样的观察塞进一个定理里。哪里有模型，我们就对它一丝不苟；哪里是课堂岔进了真实经济的肌理，我们就把这段岔话保留下来，单独装进一个注记框里。请把那些框子读作本学科的一部分，而不是没收拾干净的线头。如今，宏观经济学的问题、方法与性情都已在视野之内，下一章我们就转向其中任何一项都首先需要的那项工作：度量总量经济。
